\item Un globo (su envoltura tiene una masa de $m_b = 0.250 \text{ kg}$) se llena con helio y se amarra a una cuerda uniforme de longitud $\ell = 2.00 \text{ m}$ y masa $m = 0.050 \text{ kg}$. El globo es esférico, con un radio de $r = 0.400 \text{ m}$. Cuando se libera en aire a temperatura de $20^\circ\text{C}$ y densidad $\rho_{\text{aire}} = 1.20 \text{ kg/m}^3$, se eleva una longitud $h$ de cuerda y después permanece en equilibrio como se muestra en la figura P14.33. Se desea encontrar la longitud de cuerda levantada por el globo. (a) Cuando el globo permanece estacionario, ¿cuál es el modelo de análisis apropiado para describirlo? (b) Escriba una ecuación de fuerzas para el globo, a partir de este modelo, en términos de la fuerza boyante $B$, el peso $F_b$ del globo, el peso $F_{\text{He}}$ del helio y el peso $F_s$ del segmento de cuerda de longitud $h$. (c) Realice una apropiada sustitución para cada una de esas fuerzas y resuelva simbólicamente para la masa $m_s$ del segmento de cuerda de longitud $h$ en términos de $m_b$, $r$, $\rho_{\text{aire}}$ y la densidad del helio $\rho_{\text{He}}$. (d) Obtenga el valor numérico de la masa $m_s$. (e) Determine la longitud $h$ numéricamente.
